\uuid{Hn9L}
\titre{Domaine de définition et lignes de niveau}
\niveau{L2}
\module{Analyse}
\chapitre{Fonctions de plusieurs variables}
\sousChapitre{Domaine de définition et courbes de niveau}
\theme{Domaine de définition, logarithme, racine carrée, courbes de niveau}
\auteur{Maxime Nguyen}
\datecreate{2026-03-24}
\organisation{amscc}
\difficulte{2}

\contenu{
\texte{
  On considère la fonction $f$ définie par
  \[
    f(x,y) = \frac{\sqrt{5y - 2x}}{\ln(x-1)}.
  \]
}
\begin{enumerate}
  \item \question{Déterminer le domaine de définition de $f$.}
    \reponse{
      L'expression $f(x,y)$ est définie si et seulement si :
      \begin{itemize}
        \item $x > 1$ (pour que $\ln(x-1)$ soit bien défini),
        \item $5y - 2x \geq 0$ (pour que $\sqrt{5y-2x}$ soit bien définie),
        \item $\ln(x-1) \neq 0$, i.e. $x \neq 2$ (pour que le quotient soit bien défini).
      \end{itemize}
      Donc :
      \[
        D_f = \left\{(x,y) \in \mathbb{R}^2 \;\middle|\; x > 1,\; 5y - 2x \geq 0,\; x \neq 2\right\}.
      \]
    }

  \item \question{Déterminer la ligne de niveau $k = 0$ de $f$.}
    \reponse{
      $f(x,y) = 0$ si et seulement si $\sqrt{5y-2x} = 0$, soit $5y = 2x$, i.e. $y = \dfrac{2}{5}x$. En tenant compte du domaine, la courbe de niveau $0$ est :
      \[
        \left\{(x,y) \in \mathbb{R}^2 \;\middle|\; y = \tfrac{2}{5}x,\; x > 1,\; x \neq 2\right\}.
      \]
      Il s'agit de l'union d'un segment et d'une demi-droite.
    }
\end{enumerate}
}
